%%% =======================================================================
%%% ISLS Extension Phase 11 v1.0.0
%%% C29: The Navigator — Spectral-Guided Pattern Space Exploration
%%% TRITON Spiral + Evolutionary Triangulation + C16 Laplacian
%%% Author: Sebastian Klemm
%%% Date: 18 March 2026
%%% Status: Implementation Specification for Claude Code
%%% =======================================================================
\documentclass[11pt,a4paper,twoside]{article}

\usepackage[utf8]{inputenc}
\usepackage{amsmath,amssymb,amsthm}
\usepackage{booktabs,longtable,tabularx}
\usepackage{enumitem}
\usepackage{geometry}
\usepackage{hyperref}
\usepackage{cleveref}
\usepackage{xcolor}
\usepackage{fancyhdr}
\usepackage{listings}
\usepackage{titlesec}

\geometry{margin=2.5cm,top=3cm,bottom=3cm,headheight=14pt}

\theoremstyle{definition}
\newtheorem{definition}{Definition}[section]
\newtheorem{requirement}{Requirement}[section]

\theoremstyle{plain}
\newtheorem{theorem}{Theorem}[section]

\theoremstyle{remark}
\newtheorem{remark}{Remark}[section]

\newcommand{\ISLS}{\textsf{ISLS}}

\definecolor{sec-color}{HTML}{1e3a5f}

\titleformat{\section}
  {\Large\bfseries\color{sec-color}}
  {\thesection}{1em}{}

\lstset{
  basicstyle=\ttfamily\small,
  breaklines=true,
  frame=single,
  backgroundcolor=\color{gray!5},
  rulecolor=\color{gray!30},
}

\pagestyle{fancy}
\fancyhf{}
\fancyhead[LE,RO]{\ISLS{} Phase 11 v1.0.0}
\fancyhead[RE,LO]{Sebastian Klemm}
\fancyfoot[C]{\thepage}

\hypersetup{
  colorlinks=true,
  linkcolor=blue!60!black,
  pdftitle={ISLS Phase 11 — The Navigator},
  pdfauthor={Sebastian Klemm},
}

\begin{document}

\begin{titlepage}
\centering
\vspace*{2.5cm}
{\Huge\bfseries \ISLS{} Extension Phase~11\\[0.3cm]
v1.0.0}\\[1.5cm]
{\Large C29: The Navigator}\\[1cm]
{\large Spectral-Guided Pattern Space Exploration\\[0.3cm]
TRITON spirals. The mesh triangulates.\\
The Laplacian reads the topology.\\
The gradient pulls the spiral toward gold.}\\[2cm]
{\large Sebastian Klemm}\\[0.5cm]
{\large 18 March 2026}\\[1.5cm]

\begin{abstract}
\noindent
This document specifies \texttt{isls-navigator} (C29), the
spectral-guided exploration engine for the \ISLS{} Pattern Space.

Three components, each real and tested:
\begin{itemize}[nosep]
  \item \textbf{TRITON} (\texttt{metatron\_triton/}): golden-angle
        spiral search with momentum. Existing Rust crate, tested,
        eigenst\"andig. Used as path dependency. Not modified.
  \item \textbf{Evolutionary Triangulation} (from SPEC-003r1):
        k-NN mesh construction with resonance weighting,
        topology guard (Betti stability), entropy-driven refinement.
        Reimplemented in Rust on ISLS types. No scipy, no gudhi.
  \item \textbf{C16 Laplacian} (\texttt{isls-topology/}): spectral
        decomposition, eigenvalues, Fiedler vector. Already in ISLS,
        12 tests. Used directly.
\end{itemize}

The integration: TRITON generates candidate points. Each evaluated
point becomes a vertex in the mesh. The mesh is triangulated using
k-NN + resonance weighting. C16 computes the Laplacian on the mesh.
The spectral gradient biases TRITON's momentum toward high-resonance
regions. The topology guard rejects mesh updates that break
structural invariants. Over time, the mesh becomes a map of the
entire Pattern Space --- dense where patterns are rich, sparse where
they are barren.

\textbf{What is NOT in this spec:}
\begin{itemize}[nosep]
  \item No port of DioniceOS MeshHolo (too intertwined with MEF types).
  \item No port of Spiral Coupling (own state management, own persistence).
  \item No external libraries for Betti numbers (computed from Euler
        characteristic on the simplicial complex --- $\le 50$ lines).
  \item No Python. Everything in Rust.
\end{itemize}

\medskip\noindent\textbf{Status:} Implementation Specification.\\
\textbf{New Crate:} C29 (\texttt{isls-navigator}).\\
\textbf{External:} \texttt{metatron\_triton} (path dependency).\\
\textbf{Internal:} C16 (\texttt{isls-topology}), C17 (\texttt{isls-store}),
  C22 (\texttt{isls-artifact-ir}), C23 (\texttt{isls-forge}).
\end{abstract}
\end{titlepage}

\tableofcontents
\newpage


%%% =====================================================================
\part{What Exists and What Must Be Built}\label{part:inventory}

\section{Honest Inventory}\label{sec:inventory}

\begin{tabularx}{\textwidth}{l c l X}
\toprule
\textbf{Component} & \textbf{Exists?} & \textbf{Where} & \textbf{Action} \\
\midrule
TRITON spiral search & \checkmark & \texttt{metatron\_triton/} &
  Path dependency. Do not modify. \\
TRITON SpectralSignature & \checkmark & \texttt{metatron\_triton/} &
  Map to ISLS FiveDState. \\
TRITON CalibrationStrategy & \checkmark & \texttt{metatron\_triton/} &
  Use for Forge auto-tuning. \\
C16 Laplacian & \checkmark & \texttt{isls-topology/} &
  Use directly. \\
C16 spectral\_gap & \checkmark & \texttt{isls-topology/} &
  Use for mesh quality metric. \\
C16 Kuramoto & \checkmark & \texttt{isls-topology/} &
  Use for cluster detection. \\
SimplexMesh (k-NN) & \texttimes & --- &
  Build. $\sim$150 lines. \\
Resonance weighting & \texttimes & --- &
  Build. $\sim$30 lines. \\
Topology guard (Betti) & \texttimes & --- &
  Build. $\sim$80 lines. \\
Entropy control & \texttimes & --- &
  Build. $\sim$40 lines. \\
Spectral gradient bias & \texttimes & --- &
  Build. $\sim$50 lines. \\
Store persistence & \texttimes & --- &
  Build. $\sim$60 lines (C17 tables). \\
\midrule
DioniceOS MeshHolo & \checkmark & DioniceOS &
  \textbf{Do NOT port.} Wrong types. \\
DioniceOS SpiralCoupling & \checkmark & DioniceOS &
  \textbf{Do NOT port.} Too intertwined. \\
\bottomrule
\end{tabularx}

\textbf{Estimated new code:} $\sim$400--500 lines of Rust.
Everything else is reuse.


%%% =====================================================================
\part{Architecture}\label{part:arch}

\section{The Exploration Loop}\label{sec:loop}

\begin{lstlisting}[title={The complete loop}]
pub struct Navigator {
    triton: TritonSearch<Box<dyn Fn(&[f64]) -> SpectralSignature>>,
    mesh: SimplexMesh,
    config: NavigatorConfig,
    history: Vec<ExplorationStep>,
}

impl Navigator {
    /// One exploration step
    pub fn step(&mut self) -> ExplorationStep {
        // 1. TRITON generates next candidate point
        let triton_result = self.triton.step();
        let point = triton_result.point.clone();
        let signature = triton_result.signature.clone();

        // 2. Add point as vertex to the mesh
        let vertex_id = self.mesh.add_vertex(&point, &signature);

        // 3. Connect to k nearest neighbors (k-NN edges)
        let new_edges = self.mesh.connect_knn(vertex_id, self.config.k);

        // 4. Weight edges by resonance
        self.mesh.weight_edges_by_resonance(&new_edges);

        // 5. Form simplices from complete subgraphs (triangles)
        let new_simplices = self.mesh.detect_simplices(&new_edges);

        // 6. Topology guard: check if Betti numbers are stable
        let betti_before = self.mesh.betti_numbers();
        self.mesh.add_simplices(&new_simplices);
        let betti_after = self.mesh.betti_numbers();

        if !self.config.allow_betti_change
            && betti_before != betti_after
            && self.mesh.vertex_count() > self.config.min_vertices_for_guard
        {
            // Reject: rollback simplices
            self.mesh.remove_simplices(&new_simplices);
        }

        // 7. Compute Laplacian on the mesh graph
        let laplacian = self.mesh.laplacian_matrix();

        // 8. Spectral gradient: Fiedler vector points toward
        //    the direction of maximum graph connectivity change
        let gradient = self.mesh.spectral_gradient(vertex_id, &laplacian);

        // 9. Feed gradient back to TRITON as momentum bias
        let reward = signature.resonance()
            / self.triton.best_signature()
                .map(|s| s.resonance())
                .unwrap_or(1.0);
        self.triton.spiral_mut().update_momentum(&gradient, reward);

        // 10. Entropy check: if local entropy is high, reduce
        //     TRITON radius (exploit). If low, increase (explore).
        let local_entropy = self.mesh.local_entropy(vertex_id);
        self.triton.spiral_mut().adapt_radius(local_entropy);

        // 11. Record step
        let step = ExplorationStep {
            index: self.history.len(),
            point,
            signature,
            vertex_id,
            edges_added: new_edges.len(),
            simplices_added: new_simplices.len(),
            spectral_gap: self.mesh.spectral_gap(&laplacian),
            local_entropy,
            betti: self.mesh.betti_numbers(),
            best_resonance: self.triton.best_signature()
                .map(|s| s.resonance()).unwrap_or(0.0),
        };
        self.history.push(step.clone());
        step
    }
}
\end{lstlisting}


%%% =====================================================================
\part{Components}\label{part:components}

\section{SimplexMesh}\label{sec:mesh}

\begin{definition}[SimplexMesh]\label{def:mesh}
\begin{lstlisting}
pub struct SimplexMesh {
    vertices: Vec<MeshVertex>,
    edges: Vec<MeshEdge>,
    simplices: Vec<Simplex>,  // triangles (3-vertex), tetrahedra (4-vertex)
}

pub struct MeshVertex {
    pub id: usize,
    pub point: Vec<f64>,          // position in parameter space
    pub signature: SpectralSignature,  // (psi, rho, omega)
    pub resonance: f64,           // psi * rho * omega
}

pub struct MeshEdge {
    pub v1: usize,
    pub v2: usize,
    pub weight: f64,              // resonance-based weight
    pub distance: f64,            // Euclidean distance in param space
}

pub struct Simplex {
    pub vertices: Vec<usize>,     // 3 for triangle, 4 for tetrahedron
}
\end{lstlisting}
\end{definition}

\begin{definition}[k-NN Edge Construction]\label{def:knn}
\begin{lstlisting}
impl SimplexMesh {
    pub fn connect_knn(&mut self, vertex_id: usize, k: usize) -> Vec<MeshEdge> {
        let point = &self.vertices[vertex_id].point;
        let mut distances: Vec<(usize, f64)> = self.vertices.iter()
            .enumerate()
            .filter(|(i, _)| *i != vertex_id)
            .map(|(i, v)| (i, euclidean_distance(point, &v.point)))
            .collect();
        distances.sort_by(|a, b| a.1.partial_cmp(&b.1).unwrap());

        let mut new_edges = Vec::new();
        for (neighbor_id, dist) in distances.into_iter().take(k) {
            if !self.edge_exists(vertex_id, neighbor_id) {
                let edge = MeshEdge {
                    v1: vertex_id,
                    v2: neighbor_id,
                    weight: 0.0,  // set by weight_edges_by_resonance
                    distance: dist,
                };
                self.edges.push(edge.clone());
                new_edges.push(edge);
            }
        }
        new_edges
    }
}
\end{lstlisting}
\end{definition}

\begin{definition}[Resonance Weighting]\label{def:resweight}
Edge weight is the geometric mean of endpoint resonances, scaled
inversely by distance:
\begin{lstlisting}
pub fn weight_edges_by_resonance(&mut self, edges: &[MeshEdge]) {
    for edge in edges {
        let r1 = self.vertices[edge.v1].resonance;
        let r2 = self.vertices[edge.v2].resonance;
        let geo_mean = (r1 * r2).sqrt();
        let weight = geo_mean / (edge.distance + 1e-10);
        // Update the stored edge
        if let Some(e) = self.edges.iter_mut()
            .find(|e| e.v1 == edge.v1 && e.v2 == edge.v2) {
            e.weight = weight;
        }
    }
}
\end{lstlisting}
This means: edges between two high-resonance points that are close
together get the highest weight. The mesh naturally ``hardens''
around good regions.
\end{definition}


\section{Topology Guard}\label{sec:topo}

\begin{definition}[Betti Numbers (Lightweight)]\label{def:betti}
For a simplicial complex in low dimension ($\le 3$):
\begin{lstlisting}
impl SimplexMesh {
    pub fn betti_numbers(&self) -> Vec<usize> {
        let v = self.vertices.len();      // 0-simplices
        let e = self.edges.len();         // 1-simplices
        let f = self.simplices.iter()     // 2-simplices (triangles)
            .filter(|s| s.vertices.len() == 3).count();
        let t = self.simplices.iter()     // 3-simplices (tetrahedra)
            .filter(|s| s.vertices.len() == 4).count();

        // Connected components (b0): use union-find on edges
        let b0 = self.connected_components();

        // Euler characteristic: chi = v - e + f - t
        let chi = v as i64 - e as i64 + f as i64 - t as i64;

        // b1 = e - v + b0 (for connected graph without higher holes)
        // This is approximate but sufficient for our purpose
        let b1 = (e as i64 - v as i64 + b0 as i64).max(0) as usize;

        // b2 = chi - b0 + b1 (from Euler formula)
        let b2 = (chi - b0 as i64 + b1 as i64).max(0) as usize;

        vec![b0, b1, b2]
    }
}
\end{lstlisting}
No gudhi. No scipy. 30 lines of Rust.
\end{definition}

\begin{definition}[Topology Guard]\label{def:topoguard}
A mesh update is accepted only if it does not change the Betti
vector beyond a tolerance (default: $\Delta b_i \le 1$ for each $i$):
\begin{lstlisting}
pub fn topology_stable(before: &[usize], after: &[usize], tolerance: usize) -> bool {
    before.iter().zip(after.iter())
        .all(|(b, a)| (*b as i64 - *a as i64).unsigned_abs() as usize <= tolerance)
}
\end{lstlisting}
This prevents the mesh from developing ``holes'' (topology changes)
due to bad edge additions. The mesh grows monotonically in
topological complexity, controlled.
\end{definition}


\section{Spectral Gradient}\label{sec:gradient}

\begin{definition}[Spectral Gradient from Fiedler Vector]\label{def:specgrad}
The Fiedler vector (second eigenvector of the Laplacian) partitions
the graph into two halves. The gradient at a vertex is the direction
from the vertex toward the center of the Fiedler-positive cluster
(the ``good'' half):
\begin{lstlisting}
impl SimplexMesh {
    pub fn spectral_gradient(
        &self, vertex_id: usize, laplacian: &Vec<Vec<f64>>
    ) -> Vec<f64> {
        // Use C16's spectral decomposition
        let eigenvalues = eigendecompose(laplacian);
        let fiedler = &eigenvalues.eigenvectors[1]; // 2nd eigenvector

        // Partition: positive Fiedler = cluster A, negative = cluster B
        let vertex_fiedler = fiedler[vertex_id];

        // Centroid of the opposite cluster (where we want to go)
        let target_cluster: Vec<usize> = fiedler.iter().enumerate()
            .filter(|(_, &f)| (f > 0.0) != (vertex_fiedler > 0.0))
            .map(|(i, _)| i)
            .collect();

        if target_cluster.is_empty() {
            return vec![0.0; self.vertices[vertex_id].point.len()];
        }

        // Direction = centroid(target) - vertex.point
        let dim = self.vertices[vertex_id].point.len();
        let mut centroid = vec![0.0; dim];
        for &tid in &target_cluster {
            for d in 0..dim {
                centroid[d] += self.vertices[tid].point[d];
            }
        }
        for d in 0..dim {
            centroid[d] /= target_cluster.len() as f64;
            centroid[d] -= self.vertices[vertex_id].point[d];
        }

        // Normalize
        let norm: f64 = centroid.iter().map(|x| x * x).sum::<f64>().sqrt();
        if norm > 1e-10 {
            centroid.iter_mut().for_each(|x| *x /= norm);
        }
        centroid
    }
}
\end{lstlisting}

This is the ``invisible string'': the Fiedler vector tells TRITON
``the other side of the graph has high-resonance points you haven't
explored --- go there.''
\end{definition}


\section{Entropy Control}\label{sec:entropy}

\begin{definition}[Local Entropy]\label{def:entropy}
Shannon entropy of the resonance distribution in a vertex's
k-neighborhood:
\begin{lstlisting}
impl SimplexMesh {
    pub fn local_entropy(&self, vertex_id: usize) -> f64 {
        let neighbors = self.neighbors_of(vertex_id);
        if neighbors.is_empty() { return 0.0; }

        let resonances: Vec<f64> = neighbors.iter()
            .map(|&n| self.vertices[n].resonance)
            .collect();
        let total: f64 = resonances.iter().sum();
        if total < 1e-10 { return 0.0; }

        let probs: Vec<f64> = resonances.iter()
            .map(|r| r / total)
            .collect();

        -probs.iter()
            .filter(|&&p| p > 0.0)
            .map(|&p| p * p.ln())
            .sum::<f64>()
    }
}
\end{lstlisting}

High entropy = many similar-quality neighbors = well-explored region
$\to$ TRITON should explore elsewhere (increase radius).

Low entropy = one dominant point surrounded by poor ones = unexplored
gradient $\to$ TRITON should exploit here (decrease radius).
\end{definition}


%%% =====================================================================
\part{TRITON Integration}\label{part:triton}

\section{Path Dependency}\label{sec:tritonpath}

\begin{lstlisting}[title={Cargo.toml}]
[dependencies]
metatron-triton = { path = "../../../metatron_triton" }
isls-types      = { path = "../isls-types" }
isls-topology   = { path = "../isls-topology" }
isls-forge      = { path = "../isls-forge" }
isls-store      = { path = "../isls-store" }
\end{lstlisting}

\begin{requirement}[TRITON Unchanged]\label{req:tritonunchanged}
The \texttt{metatron\_triton} crate is used as-is. No modifications.
If the path does not resolve, the navigator compiles with a feature
gate \texttt{\#[cfg(feature = "triton")]} and falls back to a
simple random search.
\end{requirement}

\section{Evaluator Function}\label{sec:evaluator}

The evaluator bridges TRITON to the Forge:
\begin{lstlisting}
pub fn forge_evaluator(
    forge: &mut ForgeEngine,
    base_spec: &DecisionSpec,
) -> impl Fn(&[f64]) -> SpectralSignature + '_ {
    move |params: &[f64]| {
        // Map TRITON's [0,1]^n to Forge config
        let mut config = ForgeCalibrationState::from_normalized(params);

        // Run a forge with this config
        let result = forge.forge_with_config(base_spec.clone(), &config);

        match result {
            Ok(r) => SpectralSignature::new(
                r.quality_score,    // psi: quality
                r.stability_score,  // rho: stability
                r.efficiency_score, // omega: efficiency
            ),
            Err(_) => SpectralSignature::new(0.0, 0.0, 0.0),
        }
    }
}
\end{lstlisting}


%%% =====================================================================
\part{Application: Pattern Space Exploration}\label{part:patternspace}

\section{What the Navigator Explores}\label{sec:explores}

The Navigator has two modes:

\subsection{Mode 1: Forge Config Optimization}

The parameter space is the Forge's configuration:
$[\text{pmhd\_ticks}, \text{opposition}, \text{temperature},
\text{match\_threshold}, \text{retries}]$. TRITON finds the
configuration that maximizes $\psi \times \rho \times \omega$ for
a given domain. The mesh shows which configuration regions are
``gold'' and which are ``barren.''

\subsection{Mode 2: Architecture Exploration}

The parameter space is the space of possible architectures. Each
dimension encodes an architectural decision (e.g., ``include auth
module'' = $[0,1]$, ``database type'' = $[0=sqlite, 0.5=postgres,
1=none]$, ``cache layer'' = $[0,1]$). TRITON explores combinations.
The evaluator forges each combination and measures quality. The mesh
shows which architectural regions produce working software.

\subsection{Singularity Detection}

A singularity is a vertex whose resonance is $> 2\sigma$ above its
neighborhood mean:
\begin{lstlisting}
pub fn detect_singularities(&self) -> Vec<usize> {
    self.vertices.iter().enumerate()
        .filter(|(id, v)| {
            let neighbors = self.neighbors_of(*id);
            if neighbors.len() < 3 { return false; }
            let mean: f64 = neighbors.iter()
                .map(|&n| self.vertices[n].resonance).sum::<f64>()
                / neighbors.len() as f64;
            let variance: f64 = neighbors.iter()
                .map(|&n| (self.vertices[n].resonance - mean).powi(2))
                .sum::<f64>() / neighbors.len() as f64;
            let sigma = variance.sqrt();
            v.resonance > mean + 2.0 * sigma
        })
        .map(|(id, _)| id)
        .collect()
}
\end{lstlisting}
Singularities are unexplored gold: architectures that score far
better than anything nearby. They become candidates for new
templates.


%%% =====================================================================
\part{Studio Integration}\label{part:studio}

\begin{lstlisting}[title={New Studio View: Navigator}]
+--------------------------------------------------+
|  NAVIGATOR                                         |
+--------------------------------------------------+
|  Mode: [Config Optimization v]                     |
|  Steps: 47/100  Best resonance: 0.823              |
|  [ START ] [ STEP ] [ STOP ]                       |
|                                                    |
|  +-- Mesh Visualization (2D projection) ----------+|
|  |  (canvas: vertices as dots, edges as lines,    ||
|  |   color = resonance, stars = singularities)    ||
|  |                                                 ||
|  |     *           .                               ||
|  |    / \   .---.                                  ||
|  |   .---* /     \                                 ||
|  |        \.     .                                 ||
|  |         '---'                                   ||
|  +------------------------------------------------+|
|                                                    |
|  +-- Metrics Panel -------------------------------+|
|  | Vertices: 47  Edges: 142  Simplices: 89        ||
|  | Spectral gap: 0.34  Betti: [1, 2, 0]           ||
|  | Singularities: 2                                ||
|  | Local entropy: 0.72 (exploring)                 ||
|  +------------------------------------------------+|
|                                                    |
|  +-- Singularity Feed ----------------------------+|
|  | ! Vertex 23: resonance 0.91 (2.4 sigma)        ||
|  |   Config: ticks=200 opp=0.3 temp=0.0 thr=0.7  ||
|  |   [Apply to Forge] [Save as Preset]             ||
|  +------------------------------------------------+|
+--------------------------------------------------+
\end{lstlisting}


%%% =====================================================================
\part{CLI}\label{part:cli}

\begin{lstlisting}
# Start navigator in config optimization mode
isls navigate --mode config --steps 100 --domain rust

# Start navigator in architecture exploration mode
isls navigate --mode architecture --steps 50 --template rest-api

# Show navigator status
isls navigate status

# Apply best found configuration
isls navigate apply-best

# List singularities
isls navigate singularities

# Export mesh as JSON (for visualization)
isls navigate export-mesh --output mesh.json
\end{lstlisting}


%%% =====================================================================
\part{Acceptance Tests}\label{part:tests}

\begin{longtable}{l l p{5.5cm}}
\toprule
\textbf{ID} & \textbf{Name} & \textbf{Procedure} \\
\midrule
\endfirsthead \midrule \endhead \bottomrule \endfoot
AT-NV1 & Mesh vertex add & Add 10 vertices; verify mesh has 10
  vertices. \\
AT-NV2 & k-NN edges & Add 5 vertices with k=2; verify each has
  $\le 2$ edges. \\
AT-NV3 & Resonance weight & Add 2 high-res and 2 low-res vertices;
  verify high-res edge has higher weight. \\
AT-NV4 & Simplex detection & Add 3 mutually connected vertices;
  verify 1 triangle simplex formed. \\
AT-NV5 & Betti numbers & Mesh with 4 vertices, 4 edges forming a
  cycle; verify $b_0=1, b_1=1$. \\
AT-NV6 & Topology guard & Add simplex that would change Betti;
  verify rejected when guard active. \\
AT-NV7 & Spectral gradient & Build mesh with two clusters; verify
  gradient at cluster A points toward cluster B. \\
AT-NV8 & Local entropy & Uniform neighborhood: verify high entropy.
  One dominant vertex: verify low entropy. \\
AT-NV9 & TRITON integration & Run 20 navigator steps with mock
  evaluator; verify best resonance improves. \\
AT-NV10 & Singularity detection & Inject one vertex with resonance
  3$\sigma$ above neighbors; verify detected. \\
AT-NV11 & Determinism & Run same seed twice; verify identical
  mesh and history. \\
AT-NV12 & Fallback without TRITON & Disable triton feature; verify
  navigator uses random search fallback. \\
\end{longtable}


%%% =====================================================================
\part{Handoff}\label{part:handoff}

\section{What is Built vs.\ What is Reused}

\begin{tabular}{l l r}
\toprule
\textbf{Component} & \textbf{Source} & \textbf{Est.\ Lines} \\
\midrule
SimplexMesh (struct + k-NN + edges) & New & $\sim$150 \\
Resonance weighting & New & $\sim$30 \\
Simplex detection & New & $\sim$40 \\
Betti numbers (lightweight) & New & $\sim$50 \\
Topology guard & New & $\sim$20 \\
Spectral gradient (Fiedler) & New + C16 & $\sim$50 \\
Local entropy & New & $\sim$30 \\
Singularity detection & New & $\sim$20 \\
Navigator loop + config & New & $\sim$80 \\
Evaluator bridge to Forge & New & $\sim$40 \\
\midrule
TRITON spiral + search & Reuse & 0 (path dep) \\
C16 Laplacian + spectral & Reuse & 0 (existing) \\
\midrule
\textbf{Total new code} & & $\sim$\textbf{510 lines} \\
\bottomrule
\end{tabular}

\section{Test Count}

\begin{tabular}{l r r}
\toprule
\textbf{Phase} & \textbf{New Tests} & \textbf{Cumulative} \\
\midrule
Phase 1--10 & 355 & 355 \\
Phase 11 (C29) & 12 & $\ge$ 367 \\
\bottomrule
\end{tabular}

\section{Completion Criteria}

\begin{enumerate}[nosep]
  \item SimplexMesh builds mesh from evaluated points with k-NN.
  \item Resonance weighting hardens edges around high-quality regions.
  \item Betti numbers computed without external libraries.
  \item Topology guard rejects destabilizing mesh updates.
  \item C16 Laplacian computes spectral gap and Fiedler vector on mesh.
  \item Spectral gradient biases TRITON momentum toward unexplored
        high-resonance regions.
  \item Local entropy controls TRITON's explore/exploit radius.
  \item Singularity detection flags anomalously high-scoring vertices.
  \item TRITON integrated as path dependency, unmodified.
  \item Navigator runs 100 steps and converges on a mock evaluator.
  \item Total tests $\ge 367$, zero warnings.
  \item \textbf{The spiral explores. The mesh remembers.
        The spectrum guides. The topology guards.}
\end{enumerate}

\vfill
\begin{center}
\rule{0.5\textwidth}{0.4pt}\\[0.5cm]
{\small Generated for \ISLS{} Extension Phase 11 v1.0.0.\\[0.2cm]
$\sim$510 lines of new Rust.\\
Everything else is reuse: TRITON spirals, C16 analyzes, the mesh bridges.\\[0.3cm]
TRITON is the Goldfisch.\\
The mesh is the Hamsterball that grows.\\
The Laplacian is the invisible string.\\
The topology guard is the immune system.\\
Singularities are unexplored gold.}
\end{center}

\end{document}
